%%%%%%%%%%%%%%%%
%%% num 22 %%%
%%%%%%%%%%%%%%%%

\Chapter{22}

\Verse{1}
{
	\hR{וַיִּסְעוּ בְּנֵי יִשְׂרָאֵל וַיַּחֲנוּ בְּעַרְבוֹת מוֹאָב מֵעֵבֶר לְיַרְדֵּן יְרֵחוֹ׃}
}
{
	\eR{
		And the children of Israel traveled, and they camped in
		the plains of Moab, across the Jordan from Jericho.
	}
}
{
}

\Verse{2}
{
	\hE{וַיַּרְא בָּלָק בֶּן צִפּוֹר אֵת כּל אֲשֶׁר עָשָׂה יִשְׂרָאֵל לָאֱמֹרִי׃}
}
{
	\eE{
		And Balak, son of Zippor, saw everything that Israel had
		done to the Amorite.
	}
}
{
	\fC{REF: The Balaam episode is perhaps the hardest section in the Torah in which to delineate sources. Most scholars regard this three-chapter story as a composite: first, because they think of the accounts of repeated sets of ambassadors to Balaam as a doublet; and, second, because they think there is a contradiction in the story when God tells Balaam to go with the Moabites but then is angry at him for going. I am not at all certain that these things are evidence of two sources. The several embassies to Balaam, each composed of more distinguished ambassadors, may well be the original progression of the story. And the confusion over God’s sending Balaam and then being angry at him is surprising but still understandable as a single author’s development, and it is not easily resolved by separating this section into two sources in any case. Evidence of language is a stronger marker of sources than these considerations. The vast majority of the terms and phrases here that are identifiable with a particular source are typical of E, while only three are typical of J. And there is a particular cluster of terms and phrases here that are also found in Exodus 10 (E). And the deity is referred to as God (Elohim) in narration here seven times. I therefore have marked the story as wholly E, except that I have marked those three J passages so that one can observe them and make of it what one will. The first verse of the story (“And Balak, son of Zippor, saw everything that Israel had done to the Amorite”) refers to the defeat of the Amorites, which had occurred only in J, not in E. This verse, therefore, either comes from J or else was added by RJE as a means of connecting the J story of the defeat of the Amorites and the E story of the defeat of the Moabites.}
}

\Verse{3}
{
	\hE{וַיָּגר מוֹאָב מִפְּנֵי הָעָם מְאֹד כִּי רַב הוּא וַיָּקץ מוֹאָב מִפְּנֵי בְּנֵי יִשְׂרָאֵל׃}
}
{
	\eE{
		And Moab was very fearful because of the people because it
		was numerous, and Moab felt a disgust at the children of
		Israel.
	}
}
{
%	\footnote{
%		Moab. The descendants of Lot and his daughter (Gen 19:36–37). Footnote
%		22:3. it was numerous, and Moab felt a disgust. History repeats itself.
%		Moab’s fears are expressly like those of Egypt at the time of the
%		oppression of the Israelites. A cluster of parallel wordings conveys
%		this: Israel was “numerous” (cf. Exod 1:9). Moab “felt a disgust at the
%		children of Israel”; these exact words are used with regard to the
%		Egyptians (Exod 1:12). Israel is “more powerful than I” (Num 22:6; cf.
%		Exod 1:9). Israel “has covered the eye of the land”; these words are
%		used to describe the locust plague (Exod 10:5,15). Egypt set out to
%		control the too-numerous Israelites; Moab sets out to drive them away.
%		But the irony in both cases is that they perceive Israel to be a threat
%		because it is so big, but Israel is in fact no threat at all. Israel is
%		not said to have any hostile intentions at all. It is the actions that
%		Egypt and Moab take on the basis of their mistaken perceptions that
%		bring catastrophes on them. It is a strange human characteristic, taking
%		the very existence of someone who is different and strong to be a
%		threat.
%	}
}

\Verse{4}
{
	\hE{וַיֹּאמֶר מוֹאָב אֶל זִקְנֵי מִדְיָן עַתָּה יְלַחֲכוּ הַקָּהָל אֶת כּל סְבִיבֹתֵינוּ כִּלְחֹךְ הַשּׁוֹר אֵת יֶרֶק הַשָּׂדֶה וּבָלָק בֶּן צִפּוֹר מֶלֶךְ לְמוֹאָב בָּעֵת הַהִוא׃}
}
{
	\eE{
		And Moab said to the elders of Midian, “Now the community
		will lick up all of our surroundings the way an ox licks
		up the plants of a field!” And Balak, son of Zippor, was
		king of Moab at that time.
	}
}
{
%	\footnote{
%		Moab … Midian. The Moabites, who are descendants of Lot, ally with the
%		Midianites, who are direct descendants of Abraham himself (Gen 25:2).
%		The forces that caused the ancestors to separate from one another now
%		escalate into hostility among the descendants. And the kindness and
%		family fidelity that Abraham showed when he rescued Lot (Genesis 14) are
%		no longer factors in the minds of Lot’s descendants.
%	}
}

\Verse{5}
{
	\hE{וַיִּשְׁלַח מַלְאָכִים אֶל בִּלְעָם בֶּן בְּעֹר פְּתוֹרָה אֲשֶׁר עַל הַנָּהָר אֶרֶץ בְּנֵי עַמּוֹ לִקְרֹא לוֹ לֵאמֹר הִנֵּה עַם יָצָא מִמִּצְרַיִם הִנֵּה כִסָּה אֶת עֵין הָאָרֶץ וְהוּא יֹשֵׁב מִמֻּלִי׃}
}
{
	\eE{
		And he sent messengers to Balaam, son of Beor, at Pethor,
		which is on the river—the land of the children of his
		people—to call him, saying, “Here, a people has come out
		from Egypt. Here, it has covered the eye of the land. And
		it’s sitting across from me!
	}
}
{
%	\footnote{
%		Balaam, son of Beor. Balaam, son of Beor, is described in a plaster
%		inscription from Deir ‘Alla that was discovered in 1967. He is the
%		earliest person mentioned in the Bible who is also mentioned in an
%		archaeological source. 22:5 covered the eye of the land. Meaning: the
%		Israelites are so many that they make the ground invisible. This is the
%		same expression that is used to describe the plague of the locust swarm
%		in Egypt.
%	}
}

\Verse{6}
{
	\hE{וְעַתָּה לְכָה נָּא אָרָה לִּי אֶת הָעָם הַזֶּה כִּי עָצוּם הוּא מִמֶּנִּי אוּלַי אוּכַל נַכֶּה בּוֹ וַאֲגָרְשֶׁנּוּ מִן הָאָרֶץ כִּי יָדַעְתִּי אֵת אֲשֶׁר תְּבָרֵךְ מְבֹרָךְ וַאֲשֶׁר תָּאֹר יוּאָר׃}
}
{
	\eE{
		And now, go, curse this people for me, because it’s more
		powerful than I—maybe I’ll be able: we’ll strike it, and
		I’ll drive it out from the land—because I know that
		whoever you’ll bless will be blessed, and whoever you’ll
		curse will be cursed.”
	}
}
{
}

\Verse{7}
{
	\hE{וַיֵּלְכוּ זִקְנֵי מוֹאָב וְזִקְנֵי מִדְיָן וּקְסָמִים בְּיָדָם וַיָּבֹאוּ אֶל בִּלְעָם וַיְדַבְּרוּ אֵלָיו דִּבְרֵי בָלָק׃}
}
{
	\eE{
		And Moab’s elders and Midian’s elders went, and divination
		implements were in their hand, and they came to Balaam and
		spoke Balak’s words to him.
	}
}
{
}

\Verse{8}
{
	\hE{וַיֹּאמֶר אֲלֵיהֶם לִינוּ פֹה הַלַּיְלָה וַהֲשִׁבֹתִי אֶתְכֶם דָּבָר כַּאֲשֶׁר יְדַבֵּר יְהֹוָה אֵלָי וַיֵּשְׁבוּ שָׂרֵי מוֹאָב עִם בִּלְעָם׃}
}
{
	\eE{
		And he said to them, “Spend the night here tonight, and
		I’ll bring back word to you as YHWH will speak to me.” And
		the chiefs of Moab stayed with Balaam.
	}
}
{
}

\Verse{9}
{
	\hE{וַיָּבֹא אֱלֹהִים אֶל בִּלְעָם וַיֹּאמֶר מִי הָאֲנָשִׁים הָאֵלֶּה עִמָּךְ׃}
}
{
	\eE{
		And {\God} came to Balaam and said, “Who are these people
		with you?”
	}
}
{
%	\footnote{
%		God came to Balaam. It is surely significant that the Torah, a work that
%		comes from Israel, pictures the creator as communicating with a
%		non-Israelite prophet as well. It is a reminder that the Torah begins
%		with the story of the connection between God and all of the earth. The
%		narrowing of focus to Abraham and to his descendants is said to be so
%		that they will be a blessing to all the families of the earth. Indeed
%		the connection between Balaam and that promise given in God’s first
%		words to Abraham is explicit. The full wording to Abraham is: “And I’ll
%		bless those who bless you, and those who affront you I’ll curse. And all
%		the families of the earth will be blessed through you” (Gen 12:3). And
%		soon Balaam will bless Israel with the parallel words: “Those who bless
%		you: he’s blessed. And those who curse you: he’s cursed” (Num 24:9).
%	}
}

\Verse{10}
{
	\hE{וַיֹּאמֶר בִּלְעָם אֶל הָאֱלֹהִים בָּלָק בֶּן צִפֹּר מֶלֶךְ מוֹאָב שָׁלַח אֵלָי׃}
}
{
	\eE{
		And Balaam said to {\God}, “Balak, son of Zippor, king of
		Moab, sent to me:
	}
}
{
}

\Verse{11}
{
	\hE{הִנֵּה הָעָם הַיֹּצֵא מִמִּצְרַיִם וַיְכַס אֶת עֵין הָאָרֶץ עַתָּה לְכָה קָבָה לִּי אֹתוֹ אוּלַי אוּכַל לְהִלָּחֶם בּוֹ וְגֵרַשְׁתִּיו׃}
}
{
	\eE{
		‘Here is a people who came out from Egypt and has covered
		the eye of the land. Now go, execrate it for me. Maybe
		I’ll be able to fight against it, and I’ll drive it out.’”
	}
}
{
}

\Verse{12}
{
	\hE{וַיֹּאמֶר אֱלֹהִים אֶל בִּלְעָם לֹא תֵלֵךְ עִמָּהֶם לֹא תָאֹר אֶת הָעָם כִּי בָרוּךְ הוּא׃}
}
{
	\eE{
		And {\God} said to Balaam, “You shall not go with them. You
		shall not curse the people, because it is blessed.”
	}
}
{
}

\Verse{13}
{
	\hE{וַיָּקם בִּלְעָם בַּבֹּקֶר וַיֹּאמֶר אֶל שָׂרֵי בָלָק לְכוּ אֶל אַרְצְכֶם כִּי מֵאֵן יְהֹוָה לְתִתִּי לַהֲלֹךְ עִמָּכֶם׃}
}
{
	\eE{
		And Balaam got up in the morning and said to Balak’s
		chiefs, “Go to your land, because YHWH refused to allow me
		to go with you.”
	}
}
{
}

\Verse{14}
{
	\hE{וַיָּקוּמוּ שָׂרֵי מוֹאָב וַיָּבֹאוּ אֶל בָּלָק וַיֹּאמְרוּ מֵאֵן בִּלְעָם הֲלֹךְ עִמָּנוּ׃}
}
{
	\eE{
		And the chiefs of Moab got up and came to Balak and said,
		“Balaam refused to go with us.”
	}
}
{
%	\footnote{
%		Balaam refused. Like the Egyptians toward Joseph, and like the
%		Israelites toward Moses, the Moabites focus on the human rather than on
%		God. What Balaam said was, “YHWH refused to allow me to go with you.”
%		But the Moabite emissaries report, “Balaam refused to go with us.”
%		Relating to the prophet rather than to the God who communicates through
%		the prophet is presented as a fundamental human characteristic.
%	}
}

\Verse{15}
{
	\hE{וַיֹּסֶף עוֹד בָּלָק שְׁלֹחַ שָׂרִים רַבִּים וְנִכְבָּדִים מֵאֵלֶּה׃}
}
{
	\eE{
		And Balak went on again to send more numerous and
		prestigious chiefs than these.
	}
}
{
}

\Verse{16}
{
	\hE{וַיָּבֹאוּ אֶל בִּלְעָם וַיֹּאמְרוּ לוֹ כֹּה אָמַר בָּלָק בֶּן צִפּוֹר אַל נָא תִמָּנַע מֵהֲלֹךְ אֵלָי׃}
}
{
	\eE{
		And they came to Balaam and said to him, “Balak, son of
		Zippor, said this: ‘Don’t be held back from coming to me,
	}
}
{
}

\Verse{17}
{
	\hE{כִּי כַבֵּד אֲכַבֶּדְךָ מְאֹד וְכֹל אֲשֶׁר תֹּאמַר אֵלַי אֶעֱשֶׂה וּלְכָה נָּא קָבָה לִּי אֵת הָעָם הַזֶּה׃}
}
{
	\eE{
		because I’ll honor you very much, and I’ll do everything
		that you’ll say to me. And go, execrate this people for
		me.’”
	}
}
{
}

\Verse{18}
{
	\hE{וַיַּעַן בִּלְעָם וַיֹּאמֶר אֶל עַבְדֵי בָלָק אִם יִתֶּן לִי בָלָק מְלֹא בֵיתוֹ כֶּסֶף וְזָהָב לֹא אוּכַל לַעֲבֹר אֶת פִּי יְהֹוָה אֱלֹהָי לַעֲשׂוֹת קְטַנָּה אוֹ גְדוֹלָה׃}
}
{
	\eE{
		And Balaam answered, and he said to Balak’s servants, “If
		Balak would give me a houseful of silver and gold of his I
		wouldn’t be able to go against the mouth of YHWH, my {\God},
		to do something small or big.
	}
}
{
}

\Verse{19}
{
	\hE{וְעַתָּה שְׁבוּ נָא בָזֶה גַּם אַתֶּם הַלָּיְלָה וְאֵדְעָה מַה יֹּסֵף יְהֹוָה דַּבֵּר עִמִּי׃}
}
{
	\eE{
		And now, stay here, you as well, tonight, so I may know
		what YHWH will add to speak with me.”
	}
}
{
}

\Verse{20}
{
	\hE{וַיָּבֹא אֱלֹהִים אֶל בִּלְעָם לַיְלָה וַיֹּאמֶר לוֹ אִם לִקְרֹא לְךָ בָּאוּ הָאֲנָשִׁים קוּם לֵךְ אִתָּם וְאַךְ אֶת הַדָּבָר אֲשֶׁר אֲדַבֵּר אֵלֶיךָ אֹתוֹ תַעֲשֶׂה׃}
}
{
	\eE{
		And {\God} came to Balaam at night and said to him, “If these
		people have come to call you, get up, go with them. And
		just the thing that I shall speak to you—you shall do
		that.”
	}
}
{
}

\Verse{21}
{
	\hE{וַיָּקם בִּלְעָם בַּבֹּקֶר וַיַּחֲבֹשׁ אֶת אֲתֹנוֹ וַיֵּלֶךְ עִם שָׂרֵי מוֹאָב׃}
}
{
	\eE{
		And Balaam got up in the morning and harnessed his ass and
		went with the chiefs of Moab.
	}
}
{
}

\Verse{22}
{
	\hE{וַיִּחַר אַף אֱלֹהִים כִּי הוֹלֵךְ הוּא וַיִּתְיַצֵּב מַלְאַךְ יְהֹוָה בַּדֶּרֶךְ לְשָׂטָן לוֹ וְהוּא רֹכֵב עַל אֲתֹנוֹ וּשְׁנֵי נְעָרָיו עִמּוֹ׃}
}
{
	\eE{
		And {\God}’s anger flared because he was going. And an angel
		of YHWH stood up in the road as an adversary to him. And
		he was riding on his ass, and his two boys were with him.
	}
}
{
%	\footnote{
%		God’s anger flared because he was going. But did God not just tell him
%		to go (v. 20)?! The answer may lie in v. 32; see below. Footnote 22:22.
%		an angel … his ass … his two boys. Note the parallels to the story of
%		the binding of Isaac: he takes an ass and two servant boys, and he
%		encounters an angel. And a life is spared. And altars are built. And a
%		ram is sacrificed as a burnt offering. And both stories involve the
%		doubled, emphatic form of the verb “to bless” (Gen 22:17; Num 24:10).
%		And both involve the number three. And both end with the same three
%		verbs: got up, went, returned. This may thus be another signal that the
%		merit of the patriarchs protects their descendants, the children of
%		Israel. On the basis of Abraham’s obedience at Moriah, God promises (Gen
%		22:17) to make his descendants numerous, which is the thing that now
%		distresses Balak (22:3); and God promises to bless his descendants,
%		which He now does through Balaam (24:10).
%	}
}

\Verse{23}
{
	\hE{וַתֵּרֶא הָאָתוֹן אֶת מַלְאַךְ יְהֹוָה נִצָּב בַּדֶּרֶךְ וְחַרְבּוֹ שְׁלוּפָה בְּיָדוֹ וַתֵּט הָאָתוֹן מִן הַדֶּרֶךְ וַתֵּלֶךְ בַּשָּׂדֶה וַיַּךְ בִּלְעָם אֶת הָאָתוֹן לְהַטֹּתָהּ הַדָּרֶךְ׃}
}
{
	\eE{
		And the ass saw the angel of YHWH standing up in the road
		and his sword drawn in his hand, and the ass turned from
		the road and went into the field. And Balaam hit the ass
		to turn her back to the road.
	}
}
{
}

\Verse{24}
{
	\hE{וַיַּעֲמֹד מַלְאַךְ יְהֹוָה בְּמִשְׁעוֹל הַכְּרָמִים גָּדֵר מִזֶּה וְגָדֵר מִזֶּה׃}
}
{
	\eE{
		And the angel of YHWH stood in the pathway of the
		vineyards, a fence on this side and a fence on that side.
	}
}
{
}

\Verse{25}
{
	\hE{וַתֵּרֶא הָאָתוֹן אֶת מַלְאַךְ יְהֹוָה וַתִּלָּחֵץ אֶל הַקִּיר וַתִּלְחַץ אֶת רֶגֶל בִּלְעָם אֶל הַקִּיר וַיֹּסֶף לְהַכֹּתָהּ׃}
}
{
	\eE{
		And the ass saw the angel of YHWH, and she was pressed
		against the wall, and she pressed Balaam’s foot against
		the wall, and he continued hitting her.
	}
}
{
}

\Verse{26}
{
	\hE{וַיּוֹסֶף מַלְאַךְ יְהֹוָה עֲבוֹר וַיַּעֲמֹד בְּמָקוֹם צָר אֲשֶׁר אֵין דֶּרֶךְ לִנְטוֹת יָמִין וּשְׂמֹאול׃}
}
{
	\eE{
		And the angel continued passing on and stood in a narrow
		place where there was no way to turn right or left.
	}
}
{
}

\Verse{27}
{
	\hE{וַתֵּרֶא הָאָתוֹן אֶת מַלְאַךְ יְהֹוָה וַתִּרְבַּץ תַּחַת בִּלְעָם וַיִּחַר אַף בִּלְעָם וַיַּךְ אֶת הָאָתוֹן בַּמַּקֵּל׃}
}
{
	\eE{
		And the ass saw the angel of YHWH, and she lay down under
		Balaam. And Balaam’s anger flared, and he struck the ass
		with a stick.
	}
}
{
}

\Verse{28}
{
	\hE{וַיִּפְתַּח יְהֹוָה אֶת פִּי הָאָתוֹן וַתֹּאמֶר לְבִלְעָם מֶה עָשִׂיתִי לְךָ כִּי הִכִּיתַנִי זֶה שָׁלֹשׁ רְגָלִים׃}
}
{
	\eE{
		And YHWH opened the ass’s mouth, and she said to Balaam,
		“What have I done to you that you’ve struck me these three
		times?!”
	}
}
{
}

\Verse{29}
{
	\hE{וַיֹּאמֶר בִּלְעָם לָאָתוֹן כִּי הִתְעַלַּלְתְּ בִּי לוּ יֶשׁ חֶרֶב בְּיָדִי כִּי עַתָּה הֲרַגְתִּיךְ׃}
}
{
	\eE{
		And Balaam said to the ass, “Because you abused me! If
		there had been a sword in my hand I would have killed you
		by now!”
	}
}
{
%	\footnote{
%		Balaam said to the ass. As usual, the text does not provide details
%		about emotions, and so we do not know if Balaam’s answer is expressed
%		with amazement at hearing his ass talk or if he is to be pictured as
%		just answering matter-of-factly without yet realizing that something
%		extraordinary is happening. This heightens the comic effect. Some people
%		resist seeing humor in the Bible, as if being funny is contrary to the
%		book’s dignity. Their error is not that they value the Bible too much
%		but rather that they value humor too little. Humor is such an essential
%		part of human life that the Bible would be strange and incomplete
%		without it.
%	}
}

\Verse{30}
{
	\hE{וַתֹּאמֶר הָאָתוֹן אֶל בִּלְעָם הֲלוֹא אָנֹכִי אֲתֹנְךָ אֲשֶׁר רָכַבְתָּ עָלַי מֵעוֹדְךָ עַד הַיּוֹם הַזֶּה הַהַסְכֵּן הִסְכַּנְתִּי לַעֲשׂוֹת לְךָ כֹּה וַיֹּאמֶר לֹא׃}
}
{
	\eE{
		And the ass said to Balaam, “Aren’t I your ass, on whom
		you’ve ridden? From your start until this day, have I been
		accustomed to do like this to you?” And he said, “No.”
	}
}
{
}

\Verse{31}
{
	\hE{וַיְגַל יְהֹוָה אֶת עֵינֵי בִלְעָם וַיַּרְא אֶת מַלְאַךְ יְהֹוָה נִצָּב בַּדֶּרֶךְ וְחַרְבּוֹ שְׁלֻפָה בְּיָדוֹ וַיִּקֹּד וַיִּשְׁתַּחוּ לְאַפָּיו׃}
}
{
	\eE{
		And YHWH uncovered Balaam’s eyes, and he saw the angel of
		YHWH standing up in the road, and his sword drawn in his
		hand. And he knelt and bowed to his nose.
	}
}
{
}

\Verse{32}
{
	\hE{וַיֹּאמֶר אֵלָיו מַלְאַךְ יְהֹוָה עַל מָה הִכִּיתָ אֶת אֲתֹנְךָ זֶה שָׁלוֹשׁ רְגָלִים הִנֵּה אָנֹכִי יָצָאתִי לְשָׂטָן כִּי יָרַט הַדֶּרֶךְ לְנֶגְדִּי׃}
}
{
	\eE{
		And the angel of YHWH said to him, “For what have you
		struck your ass these three times? Here, I came out as an
		adversary because the way was precipitous with regard to
		me.
	}
}
{
%	\footnote{
%		the way was precipitous with regard to me. This verse may provide the
%		answer to the question of why God was angry at Balaam for going. The
%		problem is that the meaning of the key word here (Hebrew yãra) is
%		uncertain. This is its only occurrence in the Torah. (Its only other
%		occurrence in the Tanak is in Job 16:11.) It may mean that Balaam was
%		too quick to go with the Moabite chiefs. Still, this remains one of the
%		classic unsolved enigmas of the Torah. Balaam appears to be doing just
%		what he was told by God to do, yet God is angry.
%	}
}

\Verse{33}
{
	\hE{וַתִּרְאַנִי הָאָתוֹן וַתֵּט לְפָנַי זֶה שָׁלֹשׁ רְגָלִים אוּלַי נָטְתָה מִפָּנַי כִּי עַתָּה גַּם אֹתְכָה הָרַגְתִּי וְאוֹתָהּ הֶחֱיֵיתִי׃}
}
{
	\eE{
		And the ass saw me and turned in front of me these three
		times. If she hadn’t turned from in front of me, I would
		have killed you, too, by now and kept her alive!”
	}
}
{
}

\Verse{34}
{
	\hE{וַיֹּאמֶר בִּלְעָם אֶל מַלְאַךְ יְהֹוָה חָטָאתִי כִּי לֹא יָדַעְתִּי כִּי אַתָּה נִצָּב לִקְרָאתִי בַּדָּרֶךְ וְעַתָּה אִם רַע בְּעֵינֶיךָ אָשׁוּבָה לִּי׃}
}
{
	\eE{
		And Balaam said to the angel of YHWH, “I’ve sinned,
		because I didn’t know that you were standing up toward me
		in the road. And now, if it’s bad in your eyes, let me go
		back.”
	}
}
{
}

\Verse{35}
{
	\hE{וַיֹּאמֶר מַלְאַךְ יְהֹוָה אֶל בִּלְעָם לֵךְ עִם הָאֲנָשִׁים וְאֶפֶס אֶת הַדָּבָר אֲשֶׁר אֲדַבֵּר אֵלֶיךָ אֹתוֹ תְדַבֵּר וַיֵּלֶךְ בִּלְעָם עִם שָׂרֵי בָלָק׃}
}
{
	\eE{
		And the angel of YHWH said to Balaam, “Go with the people.
		And nothing but the thing that I shall speak to you: that
		is what you shall speak.” And Balaam went with Balak’s
		chiefs.
	}
}
{
}

\Verse{36}
{
	\hE{וַיִּשְׁמַע בָּלָק כִּי בָא בִלְעָם וַיֵּצֵא לִקְרָאתוֹ אֶל עִיר מוֹאָב אֲשֶׁר עַל גְּבוּל אַרְנֹן אֲשֶׁר בִּקְצֵה הַגְּבוּל׃}
}
{
	\eE{
		And Balak heard that Balaam was coming, and he went out to
		him, to a city of Moab that was on the Arnon border, that
		was at the edge of the border.
	}
}
{
}

\Verse{37}
{
	\hE{וַיֹּאמֶר בָּלָק אֶל בִּלְעָם הֲלֹא שָׁלֹחַ שָׁלַחְתִּי אֵלֶיךָ לִקְרֹא לָךְ לָמָּה לֹא הָלַכְתָּ אֵלָי הַאֻמְנָם לֹא אוּכַל כַּבְּדֶךָ׃}
}
{
	\eE{
		And Balak said to Balaam, “Didn’t I send to you, to call
		you? Why didn’t you come to me? Am I indeed not able to
		honor you?!”
	}
}
{
}

\Verse{38}
{
	\hE{וַיֹּאמֶר בִּלְעָם אֶל בָּלָק הִנֵּה בָאתִי אֵלֶיךָ עַתָּה הֲיָכֹל אוּכַל דַּבֵּר מְאוּמָה הַדָּבָר אֲשֶׁר יָשִׂים אֱלֹהִים בְּפִי אֹתוֹ אֲדַבֵּר׃}
}
{
	\eE{
		And Balaam said to Balak, “Here, I’ve come to you. Now:
		will I be able to speak anything? The thing that {\God} will
		set in my mouth: that is what I’ll speak.”
	}
}
{
}

\Verse{39}
{
	\hE{וַיֵּלֶךְ בִּלְעָם עִם בָּלָק וַיָּבֹאוּ קִרְיַת חֻצוֹת׃}
}
{
	\eE{
		And Balaam went with Balak, and they came to Kiriath
		Huzoth.
	}
}
{
}

\Verse{40}
{
	\hE{וַיִּזְבַּח בָּלָק בָּקָר וָצֹאן וַיְשַׁלַּח לְבִלְעָם וְלַשָּׂרִים אֲשֶׁר אִתּוֹ׃}
}
{
	\eE{
		And Balak sacrificed oxen and sheep, and he let them go to
		Balaam and to the chiefs who were with them.
	}
}
{
%	\footnote{
%		he let them go. He gave the sacrificed animals to Balaam and the chiefs
%		as a great meal prior to the intended cursing of Israel.
%	}
}

\Verse{41}
{
	\hE{וַיְהִי בַבֹּקֶר וַיִּקַּח בָּלָק אֶת בִּלְעָם וַיַּעֲלֵהוּ בָּמוֹת בָּעַל וַיַּרְא מִשָּׁם קְצֵה הָעָם׃}
}
{
	\eE{
		And it was in the morning, and Balak took Balaam and
		brought him up to the High Places of Baal, and he saw the
		outer edge of the people from there.
	}
}
{
%	\footnote{
%		High Places. See the comment on Num 21:28.
%	}
}
